\documentclass[a4paper, 11pt]{article}

% --- PACCHETTI ---
\usepackage[utf8]{inputenc}
\usepackage[italian]{babel}
\usepackage[margin=2.5cm]{geometry}
\usepackage{graphicx}
\usepackage{tabularx}
\usepackage{hyperref}
\usepackage{fancyhdr}
\usepackage{array}
\usepackage{titlesec}
\usepackage{xcolor}

% --- CONFIGURAZIONE HEADER E FOOTER ---
\pagestyle{fancy}
\fancyhf{}
\fancyhead[L]{Bit Crusher}
\fancyhead[R]{Verbale Interno N. x} %INSERIRE NUMERO RIUNIONE 
\fancyfoot[C]{\thepage}
\renewcommand{\headrulewidth}{0.4pt}

% Configurazione link indice
\hypersetup{
    colorlinks=true,
    linkcolor=black,
    urlcolor=blue
}

\begin{document}

% ==========================================
% FRONTESPIZIO
% ==========================================
\begin{titlepage}
    \centering
    
    % --- LOGHI ---
    % Sostituisci i \framebox con i comandi \includegraphics corretti
    \begin{minipage}{0.45\textwidth}
        \centering
        % Esempio: \includegraphics[width=0.7\textwidth]{logo_unipd.png}
        \includegraphics[width=0.7\textwidth]{unipd_logo.png} 
    \end{minipage}
    \hfill
    \begin{minipage}{0.45\textwidth}
        \centering
        % Esempio: \includegraphics[width=0.7\textwidth]{WhatsApp Image 2026-10-15 at 17.29.12.jpeg}
        \includegraphics[width=0.7\textwidth]{Logo-BitCrusher-Nero.png}
    \end{minipage}

    \vspace{1.5cm}
    {\Large \textbf{UNIVERSITÀ DEGLI STUDI DI PADOVA}} \\
    \vspace{0.3cm}
    {\large Corso di Laurea in Informatica} \\
    \vspace{0.3cm}
    {\large Insegnamento di Ingegneria del Software} \\
    \vspace{0.3cm}
    {\large Anno Accademico 2026/2027} \\

    \vspace{2cm}
    {\Huge \textbf{Verbale di Riunione}} \\
    \vspace{0.5cm}
    {\Large Riunione Interna \#x}  %INSERIRE NUMERO RIUNIONE

    \vspace{2cm}
    \begin{flushleft}
        \large
        \textbf{Gruppo:} Bit Crusher \\
        \textbf{Email:} bitcrusherteam@gmail.com \\
        \textbf{Versione:} v1.0.0 \\
        \textbf{Data della riunione:} [Inserire data riunione] \\ 
        \textbf{Data ultima modifica:} [Inserire data ultima modifica]
    \end{flushleft}
    \vfill
\end{titlepage}

\newpage

% ==========================================
% REGISTRO DELLE MODIFICHE
% ==========================================
\section*{Registro delle Modifiche}
\vspace{0.5cm}
\begin{center}
\renewcommand{\arraystretch}{1.5}
\begin{tabularx}{\textwidth}{|l|l|l|l|X|}
\hline
\textbf{Versione} & \textbf{Data} & \textbf{Autore} & \textbf{Verificatore} & \textbf{Descrizione} \\ \hline
1.0.0 & Inserire data & Nome Cognome & Nome Cognome & Inserire Descrizione \\ \hline
\end{tabularx}
\end{center}

\newpage

% ==========================================
% INDICE
% ==========================================
\tableofcontents
\newpage

% ==========================================
% 1. INFORMAZIONI GENERALI
% ==========================================
\section{Informazioni Generali}
\begin{itemize}
    \item \textbf{Luogo/Piattaforma:} Riunione Discord
    \item \textbf{Data:} [Inserire data riunione]
    \item \textbf{Orario di inizio:} [Inserire orario di inizio]
    \item \textbf{Orario di fine:} [Inserire orario di fine]
    \item \textbf{Responsabile:} [Inserire Nome]
    \item \textbf{Scriba:} [Inserire Nome]
\end{itemize}

\vspace{0.5cm}
\textbf{Partecipanti presenti:} \\
\vspace{0.2cm}
\renewcommand{\arraystretch}{1.2}
\begin{tabularx}{0.8\textwidth}{|X|c|}
\hline
\textbf{Nome e Cognome} & \textbf{Matricola} \\ \hline
Pietro Leonardo Acampora & 2146762 \\ \hline
Marco Bollettin & 2147955 \\ \hline
Erik Giuliano & 2138012 \\ \hline
Alexandru Mitu & 2101083 \\ \hline
Riccardo Pozzato & 2147988 \\ \hline
Alex Singh & 2148621 \\ \hline
Antonio Tufano & 2137972 \\ \hline
\end{tabularx}

\vspace{0.5cm}
\textbf{Partecipanti assenti:} \\
\vspace{0.2cm}
\renewcommand{\arraystretch}{1.2}
\begin{tabularx}{0.8\textwidth}{|X|c|}
\hline
\textbf{Nome e Cognome} & \textbf{Matricola} \\ \hline
Nessuno & - \\ \hline
\end{tabularx}

% ==========================================
% 2. ORDINE DEL GIORNO
% ==========================================
\section{Ordine del Giorno}
Gli argomenti previsti per la riunione odierna sono:
\begin{enumerate}
    \item 
    \item 
    \item 
\end{enumerate}

% ==========================================
% 3. RESOCONTO DELLA RIUNIONE
% ==========================================
\section{Resoconto della Riunione}

\subsection{Sezione 1}
Inserire descrizione

\subsection{Sezione 2}
Inserire descrizione

\subsection{Sezione 3}
Inserire descrizione



% ==========================================
% 4. DECISIONI PRESE
% ==========================================
\section{Decisioni Prese}
\renewcommand{\arraystretch}{1.5}
\begin{tabularx}{\textwidth}{|c|X|}
\hline
\textbf{Codice} & \textbf{Decisione} \\ \hline
V-1.1 & Inserire decisione presa \\ \hline
V-1.2 & Inserire decisione presa \\ \hline
V-1.3 & Inserire decisione presa \\ \hline
\end{tabularx}

% ==========================================
% 5. AZIONI DA SVOLGERE (TO-DO)
% ==========================================
\section{Azioni da Svolgere (To-Do)}
\begin{tabularx}{\textwidth}{|c|X|l|l|}
\hline
\textbf{Codice Rif.} & \textbf{Descrizione Task} & \textbf{Assegnatario} & \textbf{Scadenza} \\ \hline
VI-1.1 & Inserire descrizione & Nome Cognome & Inserire data \\ \hline
VI-1.2 & Inserire descrizione & Nome Cognome & Inserire data \\ \hline
VI-1.3 & Inserire descrizione & Nome Cognome & Inserire data \\ \hline

\end{tabularx}

% ==========================================
% 6. PROSSIMA RIUNIONE
% ==========================================
\section{Prossima Riunione}
\begin{itemize}
    \item \textbf{Data:} Inserire data
    \item \textbf{Ora:} Inserire ora
    \item \textbf{OdG preliminare:} Inserire OdG Preliminare
\end{itemize}

\end{document}